% ============================================================
% Chapter 5: Development
%   Section: News Acquisition: Discovery and Extraction
%     Subsection: The Shared Fetcher and the URL Guard
% Included from chapters/05-development/02-news-acquisition/news-acquisition.tex
% ============================================================

\subsection{The Shared Fetcher and the URL Guard}
\label{subsec:scraping-fetcher}

\paragraph{One fetch per page.} Every HTTP request of the acquisition
layer goes through one class, \texttt{Fetcher}
(\path{services/scraper/fetcher.py}). It sends a single \texttt{GET}
with a 20-second timeout and a User-Agent that names the system
(\texttt{InspiringNewsBot/1.0}), and returns the page as the server
sent it. The HTML parsers of the extraction cascade share that page:
before this, each strategy fetched for itself, so falling back to a
second parser doubled the requests sent to exactly the sites that were
already the hardest to scrape. Only a strategy that needs something
plain HTML cannot give, a rendered page, goes back to the network.

\paragraph{The URL guard.} All fetching happens on the server, and the
URLs come from outside: the analysis endpoint takes them from its
caller, and evidence retrieval from whatever the search engine
returned. Without a check, the service would be an HTTP proxy into its
own network. In the development stack that network contains the graph
database and the search engine, and on a cloud host the instance
metadata endpoint is one request away. The guard
(\path{services/scraper/url_guard.py}) is therefore applied to every
URL before it is fetched, and it is deny-by-address rather than
allow-by-domain: an allow-list of news domains would defeat the purpose
of an endpoint that accepts any article the user found, while what must
never be reachable is the infrastructure, which is a question of
address space. A URL is refused when:

\begin{itemize}
    \item its scheme is not \texttt{http} or \texttt{https}, or it has
    no host;
    \item its port is one of twelve ports used by internal services and
    never by a public web page (SSH, SMTP, SMB, PostgreSQL, MySQL,
    Redis, Neo4j's two ports, Elasticsearch, Memcached, MongoDB and
    Telnet);
    \item any address its host resolves to is private, loopback,
    link-local (where cloud metadata services live), reserved,
    multicast or unspecified. The name is resolved by the guard itself
    rather than trusted, because a public name can have the address
    \texttt{127.0.0.1}, and a name that resolves to both a public and a
    private address is refused.
\end{itemize}

A host that does not resolve at all is still refused, but counted as a
connection failure rather than as a block, since a mistyped or dead
domain is not the guard protecting anything. An explicit list of
allowed hosts, empty by default, exists for staging servers on a
private network, so that nobody has to switch the guard off to reach
one.

\paragraph{Redirects.} An HTTP library that follows redirects by
itself would walk straight past the guard: a public URL can redirect
to \texttt{127.0.0.1}. The fetcher therefore follows redirects itself,
one hop at a time, checking each new location with the guard, and
gives up after five.

\paragraph{Outcomes.} A failed fetch is classified once, in the same
way for discovery and for extraction, so that a 403 on a feed and a
403 on an article are counted alike. Table~\ref{tab:fetch-outcomes}
lists the outcomes; only \texttt{ok} produced an article.

\begin{table}[ht]
\centering
\small
\caption{Outcomes of an extraction attempt.}
\label{tab:fetch-outcomes}
\begin{tabularx}{\linewidth}{@{}lX@{}}
\toprule
Outcome & Meaning \\
\midrule
\texttt{ok}                & An article was extracted. \\
\texttt{too\_short}        & Parsed, but shorter than the run's minimum
                             body length: a paywall stub, a cookie wall,
                             a gallery. \\
\texttt{no\_content}       & Fetched, but no article text was found. \\
\texttt{http\_error}       & The server answered 4xx or 5xx (the status
                             is kept). \\
\texttt{timeout}, \texttt{connection\_error} & The server did not
                             answer, or could not be reached. \\
\texttt{blocked}           & Refused by the URL guard; never sent. \\
\texttt{unavailable}       & The strategy cannot run here (no headless
                             browser installed). \\
\texttt{error}             & Anything else. \\
\bottomrule
\end{tabularx}
\end{table}
